\documentclass[11pt,letterpaper]{article}
\usepackage[T1]{fontenc}
\usepackage[utf8]{inputenc}
\usepackage{lmodern}
\usepackage[margin=1in,headheight=14pt]{geometry}
\usepackage{amsmath,amssymb,amsthm,mathtools}
\usepackage{microtype,booktabs,longtable,array,enumitem}
\usepackage{xcolor,fancyhdr,needspace}
\definecolor{ink}{HTML}{243A35}
\definecolor{accent}{HTML}{556B47}
\usepackage[colorlinks=true,linkcolor=ink,citecolor=accent,urlcolor=accent]{hyperref}
\usepackage{xurl}
\usepackage{aliascnt}
\usepackage[nameinlink,noabbrev]{cleveref}
\hypersetup{pdftitle={Gamma Functions over the Surreals},pdfauthor={Research draft prepared with ChatGPT},pdfsubject={A sharp convexity classification, flat nonuniqueness, and surcomplex phase domains}}
\pagestyle{fancy}\fancyhf{}
\fancyhead[L]{\small Gamma functions over the surreals}
\fancyhead[R]{\small\thepage}
\renewcommand{\headrulewidth}{0.3pt}
\setlength{\parskip}{0.22em}
\setlength{\emergencystretch}{3em}
\widowpenalty=10000\clubpenalty=10000\displaywidowpenalty=10000
\setlist{itemsep=0.2em,topsep=0.35em}
\setcounter{tocdepth}{2}
\numberwithin{equation}{section}
\newtheorem{theorem}{Theorem}[section]
\crefname{theorem}{Theorem}{Theorems}
\theoremstyle{plain}
\newaliascnt{lemma}{theorem}
\newtheorem{lemma}[lemma]{Lemma}
\aliascntresetthe{lemma}
\crefname{lemma}{Lemma}{Lemmas}
\Crefname{lemma}{Lemma}{Lemmas}
\newaliascnt{proposition}{theorem}
\newtheorem{proposition}[proposition]{Proposition}
\aliascntresetthe{proposition}
\crefname{proposition}{Proposition}{Propositions}
\Crefname{proposition}{Proposition}{Propositions}
\newaliascnt{corollary}{theorem}
\newtheorem{corollary}[corollary]{Corollary}
\aliascntresetthe{corollary}
\crefname{corollary}{Corollary}{Corollaries}
\Crefname{corollary}{Corollary}{Corollaries}
\theoremstyle{definition}
\newaliascnt{definition}{theorem}
\newtheorem{definition}[definition]{Definition}
\aliascntresetthe{definition}
\crefname{definition}{Definition}{Definitions}
\Crefname{definition}{Definition}{Definitions}
\newaliascnt{example}{theorem}
\newtheorem{example}[example]{Example}
\aliascntresetthe{example}
\crefname{example}{Example}{Examples}
\Crefname{example}{Example}{Examples}
\theoremstyle{remark}
\newaliascnt{remark}{theorem}
\newtheorem{remark}[remark]{Remark}
\aliascntresetthe{remark}
\crefname{remark}{Remark}{Remarks}
\Crefname{remark}{Remark}{Remarks}
\newaliascnt{warning}{theorem}
\newtheorem{warning}[warning]{Warning}
\aliascntresetthe{warning}
\crefname{warning}{Warning}{Warnings}
\Crefname{warning}{Warning}{Warnings}
\newcommand{\No}{\mathbf{No}}
\newcommand{\SC}{\mathbf{No}[i]}
\newcommand{\R}{\mathbb R}
\newcommand{\C}{\mathbb C}
\newcommand{\N}{\mathbb N}
\newcommand{\Q}{\mathbb Q}
\newcommand{\Z}{\mathbb Z}
\newcommand{\Ocal}{\mathcal O}
\newcommand{\mi}{\mathfrak m}
\newcommand{\J}{\mathbb J}
\newcommand{\Minf}{\mathfrak M_{\infty}}
\newcommand{\pp}{\operatorname{pp}}
\newcommand{\st}{\operatorname{st}}
\newcommand{\supp}{\operatorname{supp}}
\newcommand{\sgn}{\operatorname{sgn}}
\newcommand{\cis}{\operatorname{cis}}
\newcommand{\Log}{\operatorname{Log}}
\newcommand{\Imn}{\operatorname{Im}}
\newcommand{\Ren}{\operatorname{Re}}
\newcommand{\der}{\partial}
\newcommand{\abs}[1]{\lvert#1\rvert}
\newcommand{\cst}{\tfrac12\log(2\pi)}
\newcommand{\tube}{\mathcal T}
\newcommand{\Lt}{\mathcal L}
\newcommand{\Gt}{\mathcal G}
\newcommand{\repo}{\texttt{VladimirReshetnikov/Surreal}}
\newcommand{\commit}{\nolinkurl{aa846271b4dcae2c055b216126a87210292ec19b}}

\begin{document}
\hypersetup{pageanchor=false}%
\begin{titlepage}
\centering
\vspace*{0.35in}
{\large\scshape Research article\par}
\vspace{0.4in}
{\LARGE\bfseries\color{ink} Gamma Functions over the Surreals\par}
\vspace{0.22in}
{\Large A sharp convexity classification, flat nonuniqueness,\\[0.12em]
and surcomplex phase domains\par}
\vspace{0.4in}
{\large September 21, 2026\par}
\vspace{0.4in}
\begin{minipage}{0.94\textwidth}
\begin{abstract}
The classical Bohr--Mollerup characterization does not merely lose uniqueness
on the surreal positive half-line: uniqueness can fail while normalization,
strict log-convexity, every finite Gauss multiplication formula, all normalized
finite-translation germs, and every Stirling coefficient are retained.
We prove this by an explicit family and an exact coefficient classification.
Let $L_0$ be the Taylor--Stirling extension of $\log\Gamma$, let $M(x)$ be the
leading Conway monomial of a positive infinite surreal $x$, and let
$\pp(x)$ be its purely infinite part. For a class assignment
$a:\Minf\to\No$, set $h_a(x)=0$ at positive finite arguments and
$h_a(x)=a(M(x))\pp(x)$ at infinite ones. Then
\[
 L_0+h_a\text{ is convex}
 \quad\Longleftrightarrow\quad
 m\,a(m)\text{ is finite for every }m\in\Minf.
\]
In the affirmative case it is strictly convex. The proof uses exact tangent
inequalities across all relative scales, not a mean-value theorem or the
insufficient assertion that a second derivative is positive.
The choice $a(m)=\lambda e^{-m}$ gives nonzero perturbations invisible to
all algebraic orders at infinity. We also construct surcomplex log-Gamma
extensions and identify their exact finite-phase exponentiation domain inside
a finite-imaginary-part tube: at infinite real part $x$, the condition is
$y\log x$ finite. Finally, compatibility with the Berarducci--Mantova scalar
derivation eliminates every infinitesimal locally constant gauge in the
stated category.
\end{abstract}
\end{minipage}
\vfill
\begin{minipage}{0.94\textwidth}\small
\textbf{Status and provenance.}
The Taylor--Stirling construction, classical Gamma identities, restricted
analytic transfer, and the Berarducci--Mantova derivation are established inputs.
The sharp gauge classification and its combined consequences are proposed
contributions proved here. No claim is made to have resolved a named published
open problem. Their absence from all published literature is not certified;
the targeted literature and repository audit is recorded in
\cref{sec:provenance}. This draft has not been independently refereed or
verified by a proof assistant. The accompanying exact symbolic checks test
finite algebraic identities, not the full surreal theorems.
\end{minipage}
\end{titlepage}
\hypersetup{pageanchor=true}%

\tableofcontents
\clearpage

\section{The question and the results}\label{sec:introduction}

\subsection{A classical characterization confronted with new scales}
For positive real arguments, the Gamma function is uniquely characterized by
positivity, the normalization $\Gamma(1)=1$, the recurrence
$\Gamma(x+1)=x\Gamma(x)$, and convexity of $\log\Gamma$. This is the
Bohr--Mollerup theorem \cite[\S5.5(iv)]{DLMF55}. We investigate its literal
surreal analogue, with convexity required at \emph{all} positive surreal
arguments and all surreal convex-combination parameters.

The issue is not whether Gamma can be extended. Costin and Ehrlich explicitly
include Gamma and log-Gamma in their extension theory, using Taylor expansions
at finite arguments and transseries at infinite ones
\cite[\S\S8,10.1(v)]{CE}. Rather, the question here is what additional
information determines an extension, once the field contains infinitely
separated scales. Does genuine global log-convexity remove the freedom left by
a difference equation? Does imposing all finite multiplication identities help?
Can agreement to every asymptotic order still conceal a nonzero discrepancy?

The results below give a sharp negative answer within a concrete and quite
large family. Their force comes from the simultaneous requirements. A generic
periodic multiplier would usually destroy convexity or multiplication
identities; a generic infinitesimal change would usually spoil local germs.
The construction here avoids all those failures, and its convexity boundary is
proved in both directions.

\subsection{Overview of the main theorems}
Every surreal number has a Conway normal form. For positive infinite $x$,
write $M(x)$ for its leading monomial with its real coefficient removed, and
$\pp(x)$ for the sum of its purely infinite terms. Precise conventions appear
in \cref{sec:preliminaries}. The baseline $L_0$ is defined in
\cref{sec:baseline}. Given any class assignment $a$ on the positive infinite
monomials, define
\begin{equation}\label{eq:gauge-intro}
 h_a(x)=
 \begin{cases}
 0,&0<x\in\Ocal,\\
 a(M(x))\pp(x),&x>\R,
 \end{cases}
 \qquad L_a=L_0+h_a,\qquad \Gamma_a=\exp L_a.
\end{equation}
Here $x>\R$ means that $x$ exceeds every real number.

\begin{theorem}[Sharp convexity and simultaneous invariance]\label{thm:main}
For the family \eqref{eq:gauge-intro}, the following statements hold.
\begin{enumerate}[label=(\roman*)]
\item $L_a$ is convex on $\No_{>0}$ if and only if
$m a(m)\in\Ocal$ for every positive infinite monomial $m$. Under this
condition $L_a$ is strictly convex.
\item For every $a$, independently of the convexity condition,
$\Gamma_a$ agrees with $\Gamma_0$ at every positive finite argument,
satisfies $\Gamma_a(1)=1$ and $\Gamma_a(x+1)=x\Gamma_a(x)$, and obeys
\begin{equation}\label{eq:gauss-intro}
 \Gamma_a(nx)=(2\pi)^{(1-n)/2}n^{nx-1/2}
 \prod_{k=0}^{n-1}\Gamma_a(x+k/n)
\end{equation}
for every ordinary integer $n\ge1$ and every $x>0$.
\item Whenever $u$ is finite and $x,x+u>0$,
\begin{equation}\label{eq:ratio-intro}
 \frac{\Gamma_a(x+u)}{\Gamma_a(x)}
 =\frac{\Gamma_0(x+u)}{\Gamma_0(x)}.
\end{equation}
Every positive-order fine derivative of $L_a$ equals that of $L_0$.
\item The condition
$\abs{h_a(x)}\prec x^{-N}$ for every infinite $x>0$ and every
$N\in\N$ is equivalent to
$\abs{a(m)}\prec m^{-N}$ for every $m\in\Minf$ and every $N\in\N$.
\end{enumerate}
\end{theorem}

The notation $b\prec c$, for $c>0$, means that $b/c$ is infinitesimal.
In particular, part (iv) is stronger than equality of a finite number of
asymptotic terms. It supplies all of them at once. Theorem
\ref{thm:main} is proved in \cref{sec:gauges,sec:flat}, using the global tangent
gap estimates of \cref{sec:convexity}.

\begin{corollary}[Explicit strengthened nonuniqueness]\label{cor:headline}
For each real $\lambda$, set
\begin{equation}\label{eq:explicit-gauge}
 a_\lambda(m)=\lambda e^{-m}.
\end{equation}
The resulting functions $\Gamma_\lambda$ are pairwise distinct and satisfy
all the properties of \cref{thm:main}, including strict log-convexity and
flatness of their logarithmic and relative differences to every algebraic
order. At $x=\omega$,
\begin{equation}\label{eq:omega-value}
 \log\Gamma_\lambda(\omega)-\log\Gamma_0(\omega)
 =\lambda\omega e^{-\omega}.
\end{equation}
Thus normalization, recurrence, strict log-convexity, all finite Gauss
formulas, the finite Taylor extension, all normalized finite-translation
germs, and the full Stirling asymptotic expansion do not jointly characterize
a surreal Gamma function.
\end{corollary}

The surcomplex result concerns actual functions, but does not assume a
canonical exponential for arbitrary infinite imaginary arguments. On
\[
 \tube=\{x+iy:x>0,\ y\in\Ocal\},
\]
we construct locally Hahn-analytic functions $\Lt_a$ extending $L_a$.
Their canonical exponentiation using finite angles is available exactly on
\begin{equation}\label{eq:omega-domain-intro}
 \Omega=\{x+iy\in\tube:x\in\Ocal\}
 \ \cup\
 \{x+iy\in\tube:x>\R,\ y\log x\in\Ocal\}.
\end{equation}
This exact assertion and its limits are proved in
\cref{sec:surcomplex}. A different sort of compatibility restores rigidity:
\cref{sec:differential} shows that the Berarducci--Mantova scalar chain rule
forces any infinitesimal locally constant log-Gamma gauge to vanish.

\subsection{Why the main result is not a transfer tautology}
All finite real identities have suitable local Taylor extensions, but the
main theorem compares arbitrary infinite arguments, including arguments of
different leading monomials and arguments whose relative separation is
smaller than every prescribed inverse power. The coefficient assignment $a$
can vary independently on different monomials. Neither its existence nor the
sharp bound $a(m)=O(m^{-1})$ follows by replacing real variables with surreal
ones in a first-order formula about Gamma.

The new work is the compatibility between this freedom and \emph{global}
convexity. A small second derivative, or the equality $L_a''=L_0''>0$, cannot
supply that compatibility by itself. The proof measures the convexity margin
at every scale and proves that it dominates precisely the allowed gauges.

\section{Normal forms, summability, and derivative conventions}
\label{sec:preliminaries}

\subsection{Class and set conventions}
We work in a classical theory of sets and classes such as NBG. The surreal
numbers form a proper-class ordered field, and $\SC=\No[i]$ is its
complexification. Each individual normal form and every family being summed
here has set-sized support. A ``class assignment'' $a$ is a class function;
no set of all class functions is formed. Explicit choices such as
\eqref{eq:explicit-gauge} are definable without selecting one object from
each member of a proper-class family.

Every particular finite calculation, together with the supports of its
operands and the countably many series used in it, can be carried out in a
set-sized Hahn workspace. Such a workspace need not be invariant under every
operation discussed in the paper. When restricted analytic transfer is used,
one may instead choose an adequately large $\No(\varepsilon)$ with
$\varepsilon$ an ordinal epsilon-number. Here $\No(\varepsilon)$ consists
of the numbers of birthday less than $\varepsilon$; equipped with the
Taylor extensions of restricted analytic functions and the surreal
exponential, it is an elementary extension of the corresponding real
structure \cite[Prop.~4.7]{vdDE}; see also \cite[Prop.~93]{CE}.
The proofs do not regard a proper class as an ordinary topological space in
ZFC, and do not invoke a proper-class sum.

\subsection{Finite parts and monomial scales}
The standard normal-form background is due to Conway and Gonshor
\cite{Conway,Gonshor}. In our notation,
\begin{equation}\label{eq:normal-form}
 x=\sum_{\alpha\in S}r_\alpha\omega^\alpha,
 \qquad r_\alpha\in\R\setminus\{0\},
\end{equation}
where $S$ is a reverse well-ordered set of surreal exponents. We use
$\supp(x)=\{\omega^\alpha:\alpha\in S\}$ for the \emph{monomial}
support, with $\supp(0)=\varnothing$; $S$ itself is the exponent support.
In increasing-exponent formal Hahn notation $t^\gamma:=\omega^{-\gamma}$,
the term $\omega^\alpha$ has exponent $-\alpha$. This is monomial notation,
not exponentiation by a varying surreal exponent. We write
\[
 \Ocal=\{x:\abs{x}<n\text{ for some }n\in\N_{>0}\},\qquad
 \mi=\{x:\abs{x}<1/n\text{ for all }n\in\N_{>0}\}.
\]
Every $x\in\Ocal$ has a unique standard part $\st(x)\in\R$ with
$x-\st(x)\in\mi$. Define
\[
 \pp(x)=\sum_{\alpha>0}r_\alpha\omega^\alpha,
 \qquad \J=\{x:\supp(x)\subseteq\{\omega^\alpha:\alpha>0\}\}.
\]
Then $\pp:\No\to\J$ is an $\R$-linear projection, its kernel is
$\Ocal$, and $x-\pp(x)$ is finite. In particular,
\begin{equation}\label{eq:pp-translation}
 \pp(x+u)=\pp(x)\quad(u\in\Ocal).
\end{equation}
This is an exact equality, not an asymptotic estimate.

For $x>\R$, define $M(x)=\omega^{\max S}$. Its leading real coefficient is
positive, so $x/M(x)$ is finite with strictly positive standard part. Set
\[
 \Minf=\{\omega^\alpha:\alpha>0\}.
\]
For positive infinite $x$ and finite $u$,
\begin{equation}\label{eq:monomial-rules}
 M(x+u)=M(x),\qquad M(nx)=M(x),\qquad \pp(nx)=n\pp(x)
 \quad(n\in\N_{>0}).
\end{equation}
The same scaling assertions hold for any positive real scalar. Distinct
monomials have different Archimedean classes: their ratio is infinite or
infinitesimal, not a finite appreciable number.

For arbitrary surreal quantities, $f=O(g)$ means that $f/g$ is finite, and
$f\prec g$ means that $f/g$ is infinitesimal; we apply these notations with
$g>0$ or to absolute values. We write $f\sim g$ when $f/g=1+\epsilon$ with
$\epsilon\in\mi$. In estimates containing several variables, all assertions
are pointwise for the indicated surreal inputs, not statements about an
unspecified sequential limit.

\subsection{Strong sums, not topological limits}
A family of Hahn series is strongly summable when the union of its monomial
supports is reverse well ordered and each monomial occurs in only finitely
many members. Its sum is then formed coefficientwise. The positive-support
summability lemma implies the following standard fact \cite{Neumann,vdDE}:
if $\epsilon_1,\ldots,\epsilon_d$ are infinitesimal, any formal power series
with real or complex coefficients can be evaluated at these arguments as a
strong sum. Formal sums, products, and substitutions with zero constant term
then commute with evaluation.

For clarity, the reason is a support statement. Reverse the order on the
exponent group, so the supports of the infinitesimals are well-ordered
positive sets. The additive monoid generated by their finite union is well
ordered, and a fixed exponent has only finitely many representations by
finite words from those positive supports. Consequently each coefficient
in a substituted formal series receives only finitely many contributions.
This also justifies the two-variable expansions used below. No analytic
radius of convergence of the real coefficient sequence is required.

In particular, factorially growing Bernoulli coefficients cause no obstacle
to evaluating the Stirling series at $x^{-1}$ for positive infinite $x$.
This fact does \emph{not} turn its ordinary partial sums into a convergent
sequence in the fine surreal order topology. Series evaluated at surreal or
surcomplex infinitesimals below have this strong, coefficientwise meaning.
Ordinary real or complex analytic series and formal series in indeterminates
are identified separately where they are used.

\subsection{Fine coordinate derivatives versus a scalar derivation}
A fine derivative is the usual epsilon--delta derivative with surreal
positive tolerances and surreal increments. For a class function $f$ on an
open interval, $f'(x)=d$ means that for every $\epsilon>0$ there is a
$\delta>0$ such that
\[
 0<\abs{u}<\delta\quad\Longrightarrow\quad
 \left|\frac{f(x+u)-f(x)}u-d\right|<\epsilon.
\]
The complex version uses the ordered modulus on $\SC$. Class quantifiers
are interpreted in the ambient class theory. All derivative assertions below
are established on explicit local neighborhoods.

A function constant on cosets modulo $\Ocal$ is finely locally constant:
the neighborhood $\abs{u}<1$ stays in the same coset. Every fine derivative
of such a function is zero, although its values on different cosets may
change. This elementary observation is not the main theorem, but it explains
why a proof of global convexity must use more than local derivatives.

In \cref{sec:differential} the symbol $\der$ denotes instead the
Berarducci--Mantova derivation \emph{of the field of scalar surreal numbers}.
It does not differentiate a class function by varying its argument. The two
operations are related only when a chain-rule compatibility statement is
proved or explicitly imposed.

\section{The Taylor--Stirling log-Gamma function}\label{sec:baseline}

\subsection{Definition on the three positive regimes}
Let $\ell(r)=\log\Gamma(r)$ for real $r>0$, and let $\gamma$ be Euler's
constant. Put
\begin{equation}\label{eq:coefficients}
 c=\frac12\log(2\pi),\qquad
 b_j=\frac{B_{2j}}{2j(2j-1)}\quad(j\ge1),
\end{equation}
where the Bernoulli numbers have generating function
$t/(e^t-1)=\sum_{j\ge0}B_jt^j/j!$.
Define $L_0:\No_{>0}\to\No$ by
\begin{equation}\label{eq:L0-finite}
 L_0(r+\epsilon)=\sum_{k\ge0}\frac{\ell^{(k)}(r)}{k!}\epsilon^k
 \quad(r\in\R_{>0},\ \epsilon\in\mi),
\end{equation}
by
\begin{equation}\label{eq:L0-small}
 L_0(x)=-\log x-\gamma x+
 \sum_{k\ge2}\frac{(-1)^k\zeta(k)}k x^k
 \quad(0<x\in\mi),
\end{equation}
and by
\begin{equation}\label{eq:L0-infinite}
 L_0(x)=\left(x-\frac12\right)\log x-x+c+R(x),
 \qquad
 R(x)=\sum_{j\ge1}b_jx^{-(2j-1)}
 \quad(x>\R).
\end{equation}
These regimes partition the positive surreal half-line. The coefficients and
normalizations in \eqref{eq:L0-small} and \eqref{eq:L0-infinite} are the
classical ones \cite[\S\S5.7,5.11]{DLMF57,DLMF511}. We define
$\Gamma_0(x)=\exp L_0(x)$ using the usual surreal exponential.

\begin{proposition}[Well-definedness and local calculus]\label{prop:baseline-calculus}
All three definitions are well-defined strong sums. The function $L_0$ is
finely differentiable to every finite order. On each of the three regimes,
its derivatives are obtained by the displayed Taylor or formal Stirling
rules, together with $(\log x)'=1/x$. In particular, for infinite $x>0$,
\begin{align}
 L_0'(x)&=\log x-\frac1{2x}+R'(x),\label{eq:digamma}\\
 R(x)&=\frac1{12x}+O(x^{-3}),&
 R'(x)&=-\frac1{12x^2}+O(x^{-4}),\label{eq:R-bounds}\\
 L_0''(x)&=\frac1x+\frac1{2x^2}+\frac1{6x^3}+O(x^{-5}).
\end{align}
For positive infinitesimal $x$,
\begin{equation}\label{eq:small-bounds}
 L_0(x)=-\log x+O(x),\qquad L_0'(x)=-x^{-1}+O(1).
\end{equation}
\end{proposition}

\begin{proof}
The summability statement follows by substituting $\epsilon$, $x$, or
$x^{-1}$, respectively, into a formal power series with real coefficients.
All three regimes are fine-open: a sufficiently small increment leaves a
point finite appreciable, positive infinitesimal, or infinite, respectively.

Here is also the needed derivative justification, rather than just a formal
notation. In the finite appreciable case the Taylor expansion can be
re-expanded jointly in $\epsilon$ and an infinitesimal increment $u$.
The formal difference after subtracting its linear term is $u^2$ times a
strongly evaluable series; on a smaller local neighborhood that series has
a fixed finite bound. The epsilon--delta derivative follows by choosing
$\abs u$ below the requested tolerance divided by that bound. Higher
finite derivatives are obtained by repeating this argument.

For \eqref{eq:L0-small}, take $\abs u< x/2$ and expand
$\log(x+u)=\log x+\log(1+u/x)$, then shrink to make $u/x$
infinitesimal. For \eqref{eq:L0-infinite}, use
$(x+u)^{-k}=x^{-k}(1+u/x)^{-k}$ and the same logarithmic identity.
The coefficients of the remainder may now contain fixed powers of $x$ or
$x^{-1}$, but are still bounded by a fixed surreal quantity on a sufficiently
small neighborhood. An epsilon--delta choice is allowed to depend on that
quantity. The joint power series in $x^{-1}$ and $u/x$ is strongly summable,
so regrouping is legitimate. This proves the fine derivative and its
iterates, not just a coefficientwise derivative in an unrelated ring.
The explicit estimates follow from the first nonzero terms.
\end{proof}

\subsection{An exact recurrence from a formal identity}
Introduce the formal log-Laurent expression
\[
 S(X)=\left(X-\frac12\right)\log X-X+c+
 \sum_{j\ge1}b_jX^{-(2j-1)}.
\]
Translations of $\log X$ are expanded by
$\log(X+t)=\log X+\log(1+t/X)$ for real $t$.

\begin{lemma}[Formal shift identity]\label{lem:formal-shift}
As a formal expression at $X=+\infty$,
\begin{equation}\label{eq:formal-shift}
 S(X+1)-S(X)=\log X.
\end{equation}
Consequently,
\begin{equation}\label{eq:baseline-recurrence}
 L_0(x+1)-L_0(x)=\log x\qquad(x>0).
\end{equation}
\end{lemma}

\begin{proof}
One rigorous derivation of \eqref{eq:formal-shift} uses only classical
asymptotic algebra. The real function $\ell$ has the complete Stirling
expansion $S$ and satisfies $\ell(X+1)-\ell(X)=\log X$
\cite{DLMF55,DLMF511}. To determine any fixed inverse-power coefficient of
the difference, retain sufficiently many terms in the two real expansions
and expand the shift by the binomial and logarithmic formulas. The remainder
is of strictly smaller order than the coefficient under examination.
Uniqueness of real asymptotic coefficients forces that coefficient to vanish.
This holds at every finite order and proves the formal identity. Equivalently,
it follows directly from the Bernoulli generating function; the supplied
symbolic program checks the cancellation through a substantial finite order.

For infinite surreal $x$, the difference in \eqref{eq:formal-shift} is an
identity of formal power series in $x^{-1}$ after the elementary terms have
been removed. Strong evaluation therefore gives an exact equality, with no
unaccounted remainder. At finite appreciable $x$, use the local Taylor
identity for the ordinary recurrence. At positive infinitesimal $x$, formula
\eqref{eq:L0-small} is exactly the Taylor series for $\ell(1+x)-\log x$.
These cases prove \eqref{eq:baseline-recurrence} everywhere.
\end{proof}

\subsection{All finite multiplication formulas}
\begin{lemma}[No nonconstant periodic inverse-power series]\label{lem:periodic}
Let $F\in X^{-1}\R[[X^{-1}]]$ and let $t\in\R\setminus\{0\}$.
If $F(X+t)=F(X)$ formally, then $F=0$.
\end{lemma}
\begin{proof}
A nonzero $F$ has a first term $dX^{-k}$ with $d\ne0$ and $k\ge1$.
The first term of $F(X+t)-F(X)$ is $-ktdX^{-k-1}$, which is nonzero.
\end{proof}

\begin{proposition}[Exact Gauss identity for the baseline]\label{prop:gauss0}
For every ordinary integer $n\ge1$ and every positive surreal $x$,
\begin{equation}\label{eq:log-gauss}
 L_0(nx)=\sum_{k=0}^{n-1}L_0(x+k/n)
 +(nx-\tfrac12)\log n-(n-1)c.
\end{equation}
Exponentiating gives \eqref{eq:gauss-intro} for $a=0$.
\end{proposition}
\begin{proof}
For finite arguments, extend the classical identity
\cite[Eq.~5.5.6]{DLMF55} by the local Taylor rules, using
\eqref{eq:L0-small} at an infinitesimal argument. For the infinite case form
the formal defect
\[
 F_n(X)=\sum_{k=0}^{n-1}S(X+k/n)-S(nX)
 +(nX-\tfrac12)\log n-(n-1)c.
\]
Expanding only the elementary leading terms shows that all terms involving
$X\log X$, $X$, $\log X$, and the constant coefficient cancel; hence
$F_n\in X^{-1}\R[[X^{-1}]]$. Under $X\mapsto X+1/n$, the difference of
the sum is $S(X+1)-S(X)=\log X$, the difference of $-S(nX)$ is
$-\log(nX)$, and the remaining linear term contributes $\log n$.
Thus $F_n(X+1/n)=F_n(X)$. Lemma \ref{lem:periodic} gives $F_n=0$,
and strong evaluation gives \eqref{eq:log-gauss}.
\end{proof}

\begin{remark}[What a formal equality does and does not decide]
The exact definition \eqref{eq:L0-infinite} chooses the value of one
strongly evaluated formal series. Equality of \emph{asymptotic coefficients}
is a weaker condition: it permits additional terms smaller than every
$x^{-N}$. The gauges below exploit this distinction. They do not contradict
uniqueness of the strong sum of a specified family.
\end{remark}

\section{Global convexity from multiscale tangent gaps}\label{sec:convexity}

\subsection{A criterion that remains valid without completeness}
For a differentiable function $F$ on an interval in an ordered field, define
its tangent gap by
\begin{equation}\label{eq:tangent-gap}
 D_F(y,x)=F(y)-F(x)-F'(x)(y-x).
\end{equation}
We will prove positivity of this expression directly.

\begin{lemma}[Tangent support and convexity]\label{lem:tangent-convex}
Let $I$ be an interval in an ordered field and let $F:I\to K$ be finely
differentiable on its interior.
If $F$ is convex, then $D_F(y,x)\ge0$ whenever $x$ is interior and $y\in I$.
Conversely, if $D_F(y,x)>0$ for every two distinct points of an open interval,
then $F$ is strictly convex on that interval.
\end{lemma}
\begin{proof}
For the first assertion, convexity gives
\[
 F(x+t(y-x))-F(x)\le t(F(y)-F(x))\qquad(0<t<1).
\]
Dividing by $t$ and using the derivative at $x$ shows that
$F'(x)(y-x)\le F(y)-F(x)$. This step uses only the defining
fine epsilon--delta inequality: a strictly negative proposed gap would
contradict that inequality for a sufficiently small positive $t$.
No sequential convergence or supremum is needed.

For the converse let $z=(1-\theta)x+\theta y$ with $x\ne y$ and
$0<\theta<1$. Apply the strict support inequality at $z$, once with
endpoint $x$ and once with endpoint $y$. Multiply by $1-\theta$ and
$\theta$, respectively, and add. The derivative terms cancel, giving
$F(z)<(1-\theta)F(x)+\theta F(y)$.
\end{proof}

Write $D_0=D_{L_0}$. The next lemma is the analytic core of the article.
Its relative-close case is especially important: the increment need not be
an ordinary real-sized perturbation, or even a fixed inverse power of $x$.

\begin{lemma}[Infinite-argument tangent gaps]\label{lem:infinite-gaps}
Suppose $x,y>\R$ and $x\ne y$. Put $q=y/x$. Then $D_0(y,x)>0$.
More precisely:
\begin{enumerate}[label=(\roman*)]
\item If $q=1+u$ with $0\ne u\in\mi$, then
\begin{equation}\label{eq:near-gap}
 D_0(y,x)\sim\frac{(y-x)^2}{2x}.
\end{equation}
\item If $q$ is finite appreciable with $r=\st(q)\ne1$, then
\begin{equation}\label{eq:comparable-gap}
 D_0(y,x)\sim x\bigl(r\log r-r+1\bigr).
\end{equation}
\item If $q$ is positive infinite, then
\begin{equation}\label{eq:up-gap}
 D_0(y,x)\sim y\log q.
\end{equation}
\item If $q$ is positive infinitesimal, then
\begin{equation}\label{eq:down-gap}
 D_0(y,x)\sim x.
\end{equation}
\end{enumerate}
\end{lemma}

\begin{proof}
Substituting \eqref{eq:L0-infinite} and \eqref{eq:digamma} gives the exact
identity
\begin{align}
 D_0(y,x)
 &=x\Phi(q)+\frac12\Psi(q)+E_R(y,x),\label{eq:gap-identity}\\
 \Phi(q)&=q\log q-q+1,\qquad
 \Psi(q)=q-1-\log q,\label{eq:kernels}\\
 E_R(y,x)&=R(y)-R(x)-R'(x)(y-x).
\end{align}
We estimate this expression in each regime.

\emph{Relative-close arguments.}
Let $q=1+u$ with $u\in\mi\setminus\{0\}$. The formal identities
\begin{align}
 \Phi(1+u)&=\frac{u^2}{2}-\frac{u^3}{6}+\frac{u^4}{12}-\cdots,
 \label{eq:phi-series}\\
 \Psi(1+u)&=\frac{u^2}{2}-\frac{u^3}{3}+\frac{u^4}{4}-\cdots
\end{align}
are strongly evaluable. Writing $k=2j-1$, the remainder is
\begin{equation}\label{eq:R-gap-series}
 E_R(y,x)=\sum_{j\ge1}b_jx^{-k}
 \bigl((1+u)^{-k}-1+ku\bigr).
\end{equation}
Each bracket is divisible by $u^2$ as a formal power series, with quadratic
coefficient $k(k+1)/2$. The joint substitution in $x^{-1}$ and $u$ is
legitimate by positive-support summability. It gives
$E_R(y,x)=O(x^{-1}u^2)$. Hence
\[
 D_0(y,x)=\frac{xu^2}{2}\bigl(1+O(u)+O(x^{-1})\bigr),
\]
where both error terms inside the parentheses are infinitesimal.
This proves \eqref{eq:near-gap} even when $u$ lies below every
$x^{-N}$: there is no discarded error independent of $u^2$.

\emph{Comparable but not relative-close arguments.}
If $r=\st(q)>0$ and $r\ne1$, the finite Taylor extension of $\log$
gives $\Phi(q)=\Phi(r)+\epsilon$ with $\epsilon\in\mi$.
The real number $\Phi(r)$ is strictly positive. The term $\Psi(q)$ is
finite, and \eqref{eq:R-bounds} gives $E_R(y,x)=O(x^{-1})$.
Thus the first term of \eqref{eq:gap-identity} dominates and yields
\eqref{eq:comparable-gap}.

\emph{An infinitely larger second argument.}
If $q>\R$, then $\log q>\R$ and
$\Phi(q)\sim q\log q$. Also $\Psi(q)=O(q)$ and
\[
 E_R(y,x)=O(y^{-1})+O(x^{-1})+O(y/x^2).
\]
After division by $y\log q$ all these error terms are infinitesimal,
proving \eqref{eq:up-gap}.

\emph{An infinitely smaller, but still infinite, second argument.}
For $q\in\mi_{>0}$, $q\log q\in\mi$, so $\Phi(q)\sim1$.
Since $y>1$, we have
$0<\log(1/q)=\log(x/y)<\log x\prec x$.
Thus $\Psi(q)\prec x$, while $E_R(y,x)$ is infinitesimal by
\eqref{eq:R-bounds}. This proves \eqref{eq:down-gap}.

The elementary exponential comparisons just used, such as
$\log x\prec x$ at positive infinite $x$, hold for the usual surreal
exponential; they also follow from the elementary exponential-field
extension theorem \cite{vdDE}. The four cases exhaust all positive $q$.
\end{proof}

\begin{lemma}[Finite and mixed tangent gaps]\label{lem:mixed-gaps}
For two distinct positive finite surreals $x,y$, $D_0(y,x)>0$.
If exactly one of $x,y$ is infinite, $D_0(y,x)$ is positive infinite.
\end{lemma}

\begin{proof}
First suppose that both arguments are finite. Choose a positive ordinary
integer $N$ with $x,y<N$. The real function
$A_N(t)=\log\Gamma(1+t)$ is analytic on a neighborhood of $[0,N]$;
its restriction can be encoded as a restricted analytic function after a
real affine change of variables. On $0<t<N$,
\[
 L_0(t)=A_N(t)-\log t,
\]
where $A_N$ is interpreted by its Taylor extension.
The classical function $\log\Gamma$ is strictly convex, and its strict
tangent inequality on $(0,N)$ is a first-order assertion involving $A_N$,
its restricted analytic derivative, $\log$, and the field operations.
The restricted analytic exponential-field transfer theorem of
van den Dries and Ehrlich \cite{vdDE} transfers this assertion to our
$x,y$. This use of transfer is on a fixed standard bounded interval;
it does not assume a globally named Gamma function on $\No$.

Now let $x$ be finite and $y$ infinite. If $x$ is appreciable then $L_0(x)$
and $L_0'(x)$ are finite, so \eqref{eq:L0-infinite} gives
\[
 D_0(y,x)\sim y\log y>\R.
\]
If $x$ is positive infinitesimal, \eqref{eq:small-bounds} gives instead
\[
 D_0(y,x)\sim y\log y+y/x>\R.
\]
Indeed the errors are sums of $O(y)$, $O(\log y)$,
$O(\abs{\log x})$, and finite terms. They are negligible relative to the
positive displayed sum: in particular
$\abs{\log x}\prec x^{-1}\prec y/x$.
There is no assumption about the relative sizes of $y\log y$ and $y/x$.

Finally let $x$ be infinite and $y$ finite. The infinite expansion implies
\[
 xL_0'(x)-L_0(x)=x+O(\log x)\sim x,
 \qquad yL_0'(x)=O(\log x)\prec x.
\]
There is a fixed real constant $C$ such that
$\log\Gamma(t)>-C$ for every positive real $t$; for example this follows
from the positive continuous real Gamma function, its behavior at zero,
and Stirling's expansion at infinity. The same bound holds for finite
positive $y$, by the bounded-interval transfer argument above. Therefore
\[
 D_0(y,x)=L_0(y)+xL_0'(x)-L_0(x)-yL_0'(x)>x/2.
\]
This remains true if $L_0(y)$ itself is positive infinite because $y$ is
extraordinarily small. All mixed cases are covered.
\end{proof}

\begin{theorem}[Global strict log-convexity of the baseline]\label{thm:baseline-convex}
For all distinct $x,y>0$,
\begin{equation}\label{eq:strict-support0}
 L_0(y)>L_0(x)+L_0'(x)(y-x).
\end{equation}
Consequently $L_0$ is strictly convex on the entire surreal positive
half-line, with arbitrary surreal convex-combination parameters.
\end{theorem}
\begin{proof}
Combine \cref{lem:infinite-gaps,lem:mixed-gaps,lem:tangent-convex}.
\end{proof}

\subsection{The geometric meaning of the estimates}
If $y=x+\Delta$ with $\Delta\prec x$, the baseline has a convexity margin
of size $\Delta^2/x$. If $\Delta$ is finite, a gauge constant under finite
translations changes nothing. If $\Delta$ is infinite, a gauge increment
of size at most $\Delta/x$ is infinitesimal relative to that margin.
This is the source of the sharp coefficient threshold $a(M(x))=O(x^{-1})$.
At arguments with different leading scales the margin is already infinite,
so any finite gauge discrepancy is harmless. These observations are made
exact in the next section.

\section{The sharp monomial-gauge classification}\label{sec:gauges}

\subsection{Algebraic invariance and local invisibility}
For an arbitrary class assignment $a:\Minf\to\No$, define $h_a,L_a$,
and $\Gamma_a$ by \eqref{eq:gauge-intro}. No regularity of $a$ is assumed.

\begin{lemma}[Gauge identities]\label{lem:gauge-identities}
For $x,x+u>0$ and $u\in\Ocal$,
\begin{equation}\label{eq:gauge-translation}
 h_a(x+u)=h_a(x).
\end{equation}
For every positive ordinary integer $n$,
\begin{equation}\label{eq:gauge-scaling}
 h_a(nx)=n h_a(x).
\end{equation}
The function $h_a$ is finely locally constant on $\No_{>0}$; in particular
$h_a^{(k)}=0$ for every finite $k\ge1$.
\end{lemma}
\begin{proof}
A finite translation cannot change an infinite argument into a finite one or
conversely. In the finite case both sides of
\eqref{eq:gauge-translation} vanish. In the infinite case use
\eqref{eq:pp-translation} and \eqref{eq:monomial-rules}.
The same rules give \eqref{eq:gauge-scaling}; the finite case again vanishes.
For local constancy restrict to
$\abs u<\min(1,x/2)$. Repeated differentiation of a locally constant
function gives zero.
\end{proof}

\begin{proposition}[Identities shared by every gauge]\label{prop:gauge-invariance}
Every $\Gamma_a$ is positive, agrees with $\Gamma_0$ at all positive finite
arguments, is normalized at $1$, and satisfies recurrence
and every finite Gauss multiplication formula. It also satisfies
\eqref{eq:ratio-intro}, and $L_a^{(k)}=L_0^{(k)}$ for all finite $k\ge1$.
\end{proposition}
\begin{proof}
Positivity follows from the ordered exponential, and agreement and
normalization follow from $h_a=0$ on positive finite inputs.
Equation \eqref{eq:gauge-translation} with $u=1$ preserves
\eqref{eq:baseline-recurrence}. For multiplication,
\[
 \sum_{k=0}^{n-1}h_a(x+k/n)=n h_a(x)=h_a(nx),
\]
which preserves \eqref{eq:log-gauss}. Exponentiate using
$\exp(v+w)=\exp(v)\exp(w)$.
For any finite $u$, subtracting $L_a(x)$ from $L_a(x+u)$ cancels the
gauge exactly and proves the normalized ratio identity.
The derivative assertion follows from \cref{lem:gauge-identities}.
\end{proof}

\begin{warning}[Which derivatives coincide]
The equality just proved concerns the derivatives of \emph{log-Gamma}.
In general $\Gamma_a^{(k)}(x)$ does not equal $\Gamma_0^{(k)}(x)$;
on a local neighborhood it equals
$e^{h_a(x)}\Gamma_0^{(k)}(x)$. Consequently normalized derivatives such as
$\Gamma_a^{(k)}/\Gamma_a$ agree, while unnormalized values need not.
\end{warning}

\subsection{Sufficiency of the finite coefficient bound}
\begin{theorem}[Convexity under the sharp bound]\label{thm:sufficiency}
If $m a(m)\in\Ocal$ for every $m\in\Minf$, then
\begin{equation}\label{eq:gauge-strict-support}
 D_{L_a}(y,x)>0\qquad(x,y>0,\ x\ne y).
\end{equation}
In particular, $L_a$ is strictly convex.
\end{theorem}

\begin{proof}
For an infinite $x$ with $m=M(x)$,
\[
 h_a(x)=(m a(m))\frac{\pp(x)}m\in\Ocal.
\]
Thus every gauge value is finite, although there need not be a single real
bound working for all monomials. By local constancy,
\begin{equation}\label{eq:perturbed-gap}
 D_{L_a}(y,x)=D_0(y,x)+h_a(y)-h_a(x).
\end{equation}

If $x,y$ are finite, the perturbation is zero and
\cref{lem:mixed-gaps} applies. If exactly one is infinite, the baseline gap
is positive infinite while the perturbation is finite, so the conclusion
again follows.

Suppose both are infinite. If $M(x)\ne M(y)$, their ratio is infinite or
infinitesimal. By \cref{lem:infinite-gaps}, the baseline gap is positive
infinite, and a finite perturbation cannot change its sign.
If $M(x)=M(y)=m$ but $\st(y/x)\ne1$, the same conclusion follows from
\eqref{eq:comparable-gap}.

It remains to consider $M(x)=M(y)=m$ and
$\Delta=y-x\prec x$ with $\Delta\ne0$. If $\Delta$ is finite,
\eqref{eq:gauge-translation} makes the perturbation exactly zero.
If $\Delta$ is infinite, linearity of the purely infinite projection gives
\[
 h_a(y)-h_a(x)=a(m)\pp(\Delta),\qquad
 \pp(\Delta)\sim\Delta.
\]
By \eqref{eq:near-gap},
\[
 \frac{\abs{h_a(y)-h_a(x)}}{D_0(y,x)}
 \sim\frac{2\abs{a(m)}x}{\abs\Delta}\in\mi,
\]
unless $a(m)=0$, in which case the numerator vanishes outright.
Here $a(m)x$ is finite because $x/m$ and $m a(m)$ are finite, whereas
$\abs\Delta$ is infinite. Thus \eqref{eq:perturbed-gap} is strictly
positive in the last case as well. Lemma \ref{lem:tangent-convex} proves
strict convexity.
\end{proof}

\subsection{Necessity and an explicit violating tangent}
\begin{theorem}[Necessity of the finite coefficient bound]\label{thm:necessity}
If $m a(m)$ is infinite in absolute value for some $m\in\Minf$, then
$L_a$ is not convex. In fact a violating tangent can be found using two
positive infinite arguments having the same leading monomial $m$.
\end{theorem}

\begin{proof}
Fix such an $m$ and abbreviate $a=a(m)\ne0$. Both $m$ and $\abs a m$
are positive infinite. Let
\[
 b=\min(m,\abs a m),\qquad \delta=\pp(\sqrt b).
\]
Then $\delta$ is positive purely infinite, $\delta\sim\sqrt b$, and
\begin{equation}\label{eq:delta-choice}
 \delta\prec m,\qquad \delta\prec\abs a m.
\end{equation}
Set $x=m$ and $y=m-\sgn(a)\delta$. These arguments are positive infinite,
$M(y)=m$, and $y$ is purely infinite. Therefore
\[
 h_a(y)-h_a(m)=-\abs a\delta.
\]
On the other hand, \eqref{eq:near-gap} gives
\[
 D_0(y,m)\sim\frac{\delta^2}{2m}\prec\abs a\delta,
\]
since the ratio is asymptotic to $\delta/(2\abs a m)$.
Using $L_a'=L_0'$ in \eqref{eq:perturbed-gap}, we obtain
$D_{L_a}(y,m)<0$. Convex differentiable functions have nonnegative tangent
gaps by \cref{lem:tangent-convex}. Thus $L_a$ is not convex.
\end{proof}

\begin{corollary}[The exact coefficient region]\label{cor:sharp-region}
Within the entire monomialwise family \eqref{eq:gauge-intro},
\[
 \boxed{\quad L_a\text{ is convex}
 \ \Longleftrightarrow\
 \forall m\in\Minf,\ a(m)=O(m^{-1}).\quad}
\]
There are no merely convex but non-strictly-convex members of this family.
\end{corollary}

The coefficient condition is pointwise in $m$. Arbitrarily different finite
bounds at different monomials are allowed. Both positive and negative gauges
are allowed. This is not a classification of \emph{all} surreal Gamma
extensions or of all possible periodic gauges; it is an exact classification
of the explicitly stated large family.

\subsection{Examples at and beyond the boundary}
\begin{example}[A non-flat boundary gauge]\label{ex:boundary}
For a real $\lambda$, take $a(m)=\lambda/m$. Then
\[
 h_a(x)=\lambda\frac{\pp(x)}{M(x)}
\]
is finite at infinite $x$, and the theorem proves strict log-convexity.
At $x=m$ it equals $\lambda$, so unless $\lambda=0$ this gauge changes the
constant order of the logarithmic asymptotics. It demonstrates that the
allowable convexity region is much larger than the flat subfamily.
\end{example}

\begin{example}[A gauge that looks locally harmless but is not convex]
Choose one infinite monomial $m_0$, put $a(m_0)=1$, and put $a(m)=0$ at
all other monomials. Every positive-order fine derivative of $L_a$ still
coincides with that of $L_0$. Nevertheless, choosing
$x=m_0$, $y=m_0-\sqrt{m_0}$ gives
\[
 D_0(y,x)\sim\frac12,\qquad h_a(y)-h_a(x)=-\sqrt{m_0},
\]
so the tangent gap is negative. This is a concrete failure of the implication
``positive second derivative implies convexity'' for the fine calculus at
hand. It also independently checks the orientation and size of the
necessity construction.
\end{example}

\section{Flat families and the strengthened failure of uniqueness}\label{sec:flat}

\subsection{An exact flatness criterion}
\begin{definition}
An assignment $a:\Minf\to\No$ is \emph{algebraically flat} if
\[
 \abs{a(m)}\prec m^{-N}\qquad(m\in\Minf,\ N\in\N).
\]
A function $h$ on $\No_{>0}$ is algebraically flat at infinite arguments if
$\abs{h(x)}\prec x^{-N}$ for every $x>\R$ and every ordinary $N\in\N$.
\end{definition}

\begin{proposition}[Coefficient flatness equals function flatness]
\label{prop:flat-equivalence}
For the family $h_a$, coefficient flatness and function flatness are
equivalent. Under either condition $L_a$ is strictly convex and
\begin{equation}\label{eq:relative-flat}
 \frac{\Gamma_a(x)}{\Gamma_0(x)}-1\prec x^{-N}
 \qquad(x>\R,\ N\in\N).
\end{equation}
\end{proposition}
\begin{proof}
Let $m=M(x)$. Since $x/m$ is finite appreciable and
$\pp(x)/m$ is finite, the estimate
$\abs{a(m)}\prec m^{-(N+1)}$ gives
$\abs{a(m)\pp(x)}\prec x^{-N}$. This proves function flatness.
Conversely, evaluate at $x=m$, where $h_a(m)=m a(m)$, and use the
function estimate at order $N$ to obtain a coefficient estimate at order
$N+1$. Together with the weaker orders this is coefficient flatness.

In particular $m a(m)$ is infinitesimal, hence finite, so
\cref{thm:sufficiency} applies. Finally $h_a(x)$ is infinitesimal and
$\exp h_a(x)-1\sim h_a(x)$ when it is nonzero, by its strongly summable
Taylor expansion. This gives \eqref{eq:relative-flat}; the zero case is
immediate.
\end{proof}

Parts (i)--(iii) of \cref{thm:main} follow from
\cref{thm:sufficiency,thm:necessity,prop:gauge-invariance}, and part (iv)
follows from \cref{prop:flat-equivalence}. This completes its proof.

\begin{warning}[Relative flatness does not bound an absolute Gamma difference]
\label{warn:absolute-difference}
The flat quantities are $L_a-L_0$ and $\Gamma_a/\Gamma_0-1$.
Multiplication by $\Gamma_0$ need not preserve flatness. For example, with
$a_1(m)=e^{-m}$, put $h=\omega e^{-\omega}$. Then
\[
 \Delta:=\Gamma_1(\omega)-\Gamma_0(\omega)
 =\Gamma_0(\omega)(e^h-1)
 \sim\Gamma_0(\omega)\omega e^{-\omega}>0.
\]
Since $(e^h-1)/h=1+\eta$ for an infinitesimal $\eta$, taking logarithms gives
\[
 \log\Delta
 =\left(\omega+\frac12\right)\log\omega-2\omega+c+R(\omega)
   +\log(1+\eta)
 \sim\omega\log\omega.
\]
Thus $\log\Delta>N\log\omega$, and hence $\Delta>\omega^N$, for every
ordinary $N\in\N$. The absolute difference can therefore be enormous
despite agreement to every relative algebraic order.
\end{warning}

\subsection{Explicit families with arbitrarily many parameters}
\begin{proof}[Proof of \cref{cor:headline}]
For every positive infinite $m$ and ordinary $N$,
$m^N e^{-m}\in\mi$. To see this directly, the real exponential eventually
exceeds every prescribed positive real multiple of $t^{N+1}$; the same
inequality holds at infinite surreal $t$ by exponential-field transfer
\cite{vdDE}. Consequently $a_\lambda(m)=\lambda e^{-m}$ is flat for
real $\lambda$. The main theorem gives all the asserted identities and
strict convexity. At $x=\omega$, $M(x)=\pp(x)=\omega$, yielding
\eqref{eq:omega-value}. Distinct real $\lambda$ give distinct logarithms
there, and injectivity of the surreal exponential gives distinct Gamma
functions.
\end{proof}

\begin{theorem}[Independent monomial parameters]\label{thm:independent}
Let $S\subseteq\Minf$ be any set. To each function
$\lambda:S\to\R$ associate the definable class assignment
\[
 a_\lambda(m)=
 \begin{cases}
 \lambda(m)e^{-m},&m\in S,\\
 0,&m\notin S.
 \end{cases}
\]
The map $\lambda\mapsto\Gamma_{a_\lambda}$ is injective, and all its
members have the properties of \cref{cor:headline}.
\end{theorem}
\begin{proof}
Flatness is pointwise in $m$, so no boundedness of the real parameters as
$m$ varies is needed. Distinct parameter functions differ at some $m\in S$,
where their log-Gamma values differ by the nonzero quantity
$(\lambda(m)-\mu(m))m e^{-m}$.
\end{proof}

Formally, this indexed family is the single class function
$(\lambda,x)\mapsto\Gamma_{a_\lambda}(x)$ on
$\R^S\times\No_{>0}$. ``Injective'' means that different indices are
separated by some argument; it does not make proper classes into values of
a set function. This gives independent parameters on sets of any size
available in the ambient set theory, without forming a set of all class
functions.

\begin{example}[The same local data at visibly different arguments]
For the one-parameter family and $\lambda\ne0$,
\begin{align*}
 h_\lambda(\omega)&=\lambda\omega e^{-\omega},\\
 h_\lambda(\omega+17+\omega^{-1})&=\lambda\omega e^{-\omega},\\
 h_\lambda(2\omega)&=2\lambda\omega e^{-\omega},\\
 h_\lambda(\omega+\sqrt\omega)&=
 \lambda(\omega+\sqrt\omega)e^{-\omega},\\
 h_\lambda(\omega^2)&=\lambda\omega^2e^{-\omega^2}.
\end{align*}
The fourth line changes the gauge because the increment is infinite,
although it is infinitesimal relative to the argument. The second line does
not change it because the entire increment is finite. In the third line the
exponent is still $-\omega$, not $-2\omega$: $M(2\omega)=\omega$.
Confusing a leading monomial with a leading term would invalidate the
multiplication identities.
\end{example}

\subsection{Even the usual signed remainder bounds survive}
Let
\begin{equation}\label{eq:stirling-truncation}
 S_N(x)=\left(x-\frac12\right)\log x-x+c+
 \sum_{j=1}^{N}b_jx^{-(2j-1)}\qquad(N\in\N).
\end{equation}
An empty sum is zero. The following strengthens equality of asymptotic
coefficients by retaining the direction and size of each ordinary-order
Stirling remainder.

\begin{theorem}[Enveloping inequalities are still non-unique]
\label{thm:enveloping}
For every algebraically flat $a$, every $x>0$, and every ordinary $N\in\N$,
\begin{equation}\label{eq:enveloping}
 0<(-1)^N\bigl(L_a(x)-S_N(x)\bigr)
 <\abs{b_{N+1}}x^{-(2N+1)}.
\end{equation}
Thus adding all these inequalities to the requirements in
\cref{cor:headline} does not restore uniqueness.
\end{theorem}
\begin{proof}
For finite positive $x$, $L_a=L_0$ and the classical strict real remainder
bounds for log-Gamma \cite[\S5.11(ii)]{DLMF511} transfer on a standard
bounded interval, exactly as in \cref{lem:mixed-gaps}. The elementary
terms in $S_N$ introduce only powers and logarithms. For completeness, strictness follows directly from Binet's real formula
\cite[Eq.~5.9.10]{DLMF59}:
\[
 \ell(x)-S_0(x)=2\int_0^\infty
 \frac{\arctan(t/x)}{e^{2\pi t}-1}\,dt\qquad(x\in\R_{>0}).
\]
The finite expansion
\[
 \arctan v=\sum_{k=0}^{N-1}\frac{(-1)^k v^{2k+1}}{2k+1}
 +(-1)^N\int_0^v\frac{u^{2N}}{1+u^2}\,du
\]
has a signed remainder strictly between $0$ and $v^{2N+1}/(2N+1)$
for $v>0$. Integrating against the positive Binet kernel gives the
strict first-omitted-term bounds; the integrated polynomial coefficients
are the $b_j$ in \eqref{eq:coefficients}. This is an ordinary real
integral argument, used only before transfer, not a surreal integral.

For infinite $x$, the coefficients $b_j$ alternate in sign, with
$\sgn b_j=(-1)^{j+1}$. Therefore the strong-series remainder has the form
\[
 (-1)^N(L_0(x)-S_N(x))
 =\abs{b_{N+1}}x^{-(2N+1)}
  -\abs{b_{N+2}}x^{-(2N+3)}+O(x^{-(2N+5)}).
\]
Flatness makes $h_a(x)$ infinitesimal relative to the second displayed term.
The remainder itself is positive because of its first term, and the
difference from the upper bound is positive because of the second term.
This proves both strict inequalities after adding the gauge.
\end{proof}

\subsection{What ``all orders'' means here}
All differentiation orders, multiplication indices, and Stirling truncation
indices in this article are ordinary nonnegative integers. A transfinite
product, a hyperinteger-indexed recurrence scheme, or a prescription for
summing a proper-class expansion would be a different additional structure.
None is smuggled into \eqref{eq:gauss-intro} or \eqref{eq:enveloping}.

Likewise, flatness here does not mean zero. At an infinite $x$ the ordered
field contains nonzero elements smaller than every ordinary inverse power
of $x$. The explicit examples use that room. Requiring equality to a chosen
full strong sum would remove it by definition, but would be a stronger
normalization than the classical asymptotic expansion and its usual
finite-order error inequalities.

\section{Local Gamma calculus and the failure of classical propagation}
\label{sec:local}

\subsection{Normalized germs and polygamma identities}
For a positive argument $x$, the local multiplicative germ of $\Gamma_a$
normalized at $x$ is
\[
 u\longmapsto\frac{\Gamma_a(x+u)}{\Gamma_a(x)}.
\]
It is exactly the baseline germ. This includes finite translations which
are not infinitesimal, as long as positivity of the shifted argument is
retained. In particular, all logarithmic derivatives and all normalized
finite-order derivatives coincide.

Writing $\psi=L_0'$ gives, by differentiating the already established
identities,
\begin{align}
 \psi(x+1)-\psi(x)&=x^{-1},\\
 n\psi(nx)&=\sum_{k=0}^{n-1}\psi(x+k/n)+n\log n.
\end{align}
Higher fine derivatives give their polygamma counterparts. At positive
infinite $x$,
\[
 \psi(x)=\log x-\frac1{2x}-\frac1{12x^2}
 +\frac1{120x^4}-\frac1{252x^6}+\cdots,
\]
with strong summation, and
\[
 (-1)^{r+1}\psi^{(r)}(x)\sim(r-1)!x^{-r}>0
 \qquad(r\ge1).
\]
At finite positive arguments the same sign is obtained from the ordinary
polygamma functions by restricted analytic transfer. These facts agree with
the classical formulas \cite{DLMF515}, but do not distinguish any two
members of our family.

\begin{warning}[An illegitimate substitution into an ordinary sum]
For real $x>0$, the familiar identity
$\psi'(x)=\sum_{k\ge0}(x+k)^{-2}$ is an ordinarily convergent sum.
At infinite surreal $x$, its summands all have the same leading monomial
$M(x)^{-2}$, with nonzero real leading coefficient. The family is therefore
\emph{not} strongly summable: that monomial occurs infinitely often.
The correct expression used here is the differentiated Stirling strong
series, not that ordinary series with an infinite number substituted into it.
\end{warning}

\subsection{Why the real uniqueness argument does not transfer wholesale}
The unit recurrence propagates values along $x+\Z$, and finite
multiplication formulas relate a finite number of comparable arguments.
Neither operation moves an infinite argument to the real interval $(0,1]$.
More strongly, our normalized germs already prescribe every finite
translation from every starting point, and the gauge still survives.

Real Bohr--Mollerup arguments eventually use the Archimedean property to
squeeze a difference to zero through real error bounds. In the present
field, an error smaller than $x^{-N}$ for every ordinary $N$ can be nonzero.
Theorem \ref{thm:enveloping} makes the limitation explicit even for signed
Stirling inequalities. The failure is not a defect in the classical theorem:
its ambient ordered field and its globally available means of propagation
have changed.

There is no contradiction with elementary extension of the restricted
analytic exponential structure. That language names analytic functions on
fixed real compact boxes and the exponential. The normal-form projection
$\pp$, the leading-monomial selector $M$, and an arbitrary class assignment
$a$ are not silently included as definable functions in that structure.
Our use of transfer proves bounded-regime statements about $L_0$; it does
not claim definability or uniqueness of $L_a$ in an o-minimal expansion.

\subsection{Local analyticity is not global analytic continuation}
The functions $L_a$ have local Hahn-Taylor descriptions: near any point,
$h_a$ is a fixed scalar and the remaining Taylor series is the baseline one.
This does not provide a classical path joining the point to the real axis
through a connected complex domain, nor an identity theorem that propagates
an equality between different fine components. A local power-series
condition on this topology is therefore intentionally distinguished from a
single global transseries prescription.

The distinction is useful for specification as well as for theory. A
symbolic system may implement the baseline by an exact normal-form rule.
Alternatively it may specify only a recurrence, local derivatives, and
asymptotics. These are different interfaces: the second does not determine
the first, even after global log-convexity and all finite multiplication
formulas are required.

\section{Surcomplex log-Gamma and an exact finite-phase tube}
\label{sec:surcomplex}

\subsection{The canonical exponential with finite imaginary part}
For a finite surreal $b=r+\epsilon$, where $r\in\R$ and
$\epsilon\in\mi$, define
\begin{equation}\label{eq:cis}
 \cis(b)=e^{ir}\sum_{k\ge0}\frac{(i\epsilon)^k}{k!}.
\end{equation}
The sum is strong. Formal exponential identities give
$\cis(b+d)=\cis(b)\cis(d)$ for finite $b,d$ and
$\abs{\cis(b)}=1$. Thus
\begin{equation}\label{eq:finite-exp}
 E(v+ib)=e^v\cis(b),\qquad v\in\No,\ b\in\Ocal,
\end{equation}
is a canonical additive-to-multiplicative homomorphism on
$\No+i\Ocal$. It is finely complex differentiable there, with $E'=E$:
on a sufficiently small complex neighborhood the increment has a strong
exponential series and $E(w+u)=E(w)\sum u^k/k!$.

No value for $\cis(b)$ with infinite $b$ is prescribed by
\eqref{eq:cis}. The following construction keeps this domain restriction
visible. In particular, an infinite imaginary part of a log-Gamma value is
not automatically an argument at which \eqref{eq:finite-exp} may be used.

\subsection{The logarithm on the right half-plane}
For $z=x+iy$ with $x>0$, define the right-half-plane logarithm by
\[
 \Log z=\log\abs z+i\operatorname{Arg}(z),\qquad
 -\pi/2<\operatorname{Arg}(z)<\pi/2.
\]
The angle is a finite surreal. It may be constructed without evaluating a
trigonometric function at an infinite argument: use the Taylor extension of
$\arctan$ at finite $y/x$, and for infinite $\abs{y/x}$ use
$\sgn(y)\pi/2-\arctan(x/y)$. These formulas agree on their overlaps by
real analytic identities and infinitesimal Taylor evaluation. They give
$E(\Log z)=z$ and locally $\Log(z+u)=\Log z+\log(1+u/z)$, with the
last logarithm evaluated at a sufficiently small infinitesimal ratio.
Hence $\Log'(z)=1/z$.

Now restrict attention to the fine-open tube
\[
 \tube=\{x+iy\in\SC:x>0,\ y\in\Ocal\}.
\]
If $x>\R$ and $y$ is finite, a particularly simple formula is
\begin{equation}\label{eq:log-infinite-complex}
 \Log(x+iy)=\log x+\sum_{k\ge1}
 \frac{(-1)^{k+1}}k\left(\frac{iy}{x}\right)^k.
\end{equation}

\subsection{Definition and functional identities of complex log-Gamma}
Let $\ell_{\C}$ be the ordinary holomorphic branch of log-Gamma on the
complex right half-plane, chosen to equal $\log\Gamma$ on its positive
real axis. At nonzero points of its imaginary boundary, this branch has a
local analytic continuation, which is uniquely specified by continuation
from the right half-plane.

For finite $z\in\tube$ with $\zeta=\st(z)\ne0$, define
$\Lt_0(z)$ by the complex Taylor series of that local branch at $\zeta$.
Here standard part is taken componentwise. For infinitesimal $z\in\tube$,
define
\begin{equation}\label{eq:complex-small}
 \Lt_0(z)=-\Log z-\gamma z+
 \sum_{k\ge2}\frac{(-1)^k\zeta(k)}k z^k.
\end{equation}
For $z=x+iy\in\tube$ with $x>\R$, define
\begin{equation}\label{eq:complex-infinite}
 \Lt_0(z)=\left(z-\frac12\right)\Log z-z+c+
 \sum_{j\ge1}b_jz^{-(2j-1)}.
\end{equation}
The inverse $z^{-1}$ is infinitesimal, so this is a well-defined strong sum.
Finally set
\begin{equation}\label{eq:complex-gauge}
 \Lt_a(x+iy)=\Lt_0(x+iy)+h_a(x).
\end{equation}
The gauge takes values in $\No\subset\SC$, so its imaginary part is zero;
its values need not lie in the ordinary real field $\R$.

\begin{theorem}[Surcomplex logarithmic invariance]\label{thm:complex-log}
For every assignment $a$, $\Lt_a$ is finely holomorphic on $\tube$ and
has local strongly evaluable Taylor expansions. It restricts to $L_a$ on
the positive real-surreal axis, and all its positive-order fine complex
derivatives coincide with those of $\Lt_0$. Moreover,
\begin{align}
 \Lt_a(z+1)-\Lt_a(z)&=\Log z,\label{eq:complex-recurrence}\\
 \Lt_a(nz)&=\sum_{k=0}^{n-1}\Lt_a(z+k/n)
 +(nz-\tfrac12)\log n-(n-1)c\label{eq:complex-gauss}
\end{align}
for all $z\in\tube$ and ordinary $n\ge1$.
\end{theorem}

\begin{proof}
Local strong summability and differentiability follow exactly as in
\cref{prop:baseline-calculus}, with complex coefficients. At finite
standard part one uses a small ordinary complex disk on which the chosen
branch is analytic; at standard part zero use \eqref{eq:complex-small};
at infinite real part use the joint expansions in $z^{-1}$ and a small
increment divided by $z$. The extra term $h_a(\Ren z)$ is constant on
small fine disks because the real part of a sufficiently small increment
is finite. Hence it has zero fine complex derivatives.

For finite $z$, the classical analytic recurrence and multiplication
identities extend by these Taylor expansions. Their logarithmic constants
are fixed by agreement on the ordinary positive real axis. The identities
remain valid in local disks meeting a nonzero imaginary boundary point by
ordinary analytic continuation. Near zero one uses the recurrence-based
formula \eqref{eq:complex-small}.

For infinite real part, the formal proofs of
\cref{lem:formal-shift,prop:gauss0} evaluate at $z$ just as at $x$;
\eqref{eq:log-infinite-complex} ensures consistent logarithmic branches,
in particular $\Log(nz)=\log n+\Log z$. The gauge cancels because
$h_a(x+1)=h_a(x)$ and
$\sum_{k=0}^{n-1}h_a(x+k/n)=h_a(nx)$.
\end{proof}

\subsection{The exact phase estimate}
The imaginary part of \eqref{eq:complex-infinite} grows at the scale
$y\log x$, even though $y$ itself is finite. This identifies precisely when
the canonical exponential \eqref{eq:finite-exp} is available.

\begin{lemma}[Imaginary log-Gamma at infinite real part]\label{lem:phase}
Let $x>\R$ and $y\in\Ocal$. Then
\begin{equation}\label{eq:phase-estimate}
 \Imn\Lt_0(x+iy)-y\log x\in\mi.
\end{equation}
More precisely, for $y\ne0$,
\begin{equation}\label{eq:phase-quotient}
 \frac{\Imn\Lt_0(x+iy)}y
 =\log x-\frac1{2x}+O(x^{-2})+O(y^2x^{-2}).
\end{equation}
Consequently $\Imn\Lt_a(x+iy)$ is finite if and only if $y\log x$ is
finite, for every $\No$-valued gauge \eqref{eq:complex-gauge}.
\end{lemma}

\begin{proof}
Re-expand \eqref{eq:complex-infinite} around $x$ in the ratio $iy/x$.
All Taylor coefficients at $x\in\No_{>0}$ lie in $\No$, and the imaginary part
is the odd Taylor portion:
\[
 \Imn\Lt_0(x+iy)
 =yL_0'(x)-\frac{y^3}{3!}L_0'''(x)
 +\frac{y^5}{5!}L_0^{(5)}(x)-\cdots.
\]
This is an identity of strongly summable joint expansions, not an appeal to
a real convergence disk of infinite radius. Formula \eqref{eq:digamma}
gives $L_0'(x)=\log x-1/(2x)+O(x^{-2})$.
The cubic and higher odd terms, after division by $y$, have total size
$O(y^2x^{-2})$, since $L_0'''(x)=O(x^{-2})$ and each further odd
term gains a factor of order $y^2/x^2$. Strong summability justifies the
factorization. This proves \eqref{eq:phase-quotient}; if $y=0$ the
imaginary part is exactly zero. Multiplication by a finite $y$ gives
\eqref{eq:phase-estimate}. Adding $h_a(x)\in\No$ changes no imaginary
part. A quantity and an infinitesimally close quantity are finite together.
\end{proof}

\begin{theorem}[Exact canonical finite-phase Gamma domain]
\label{thm:phase-domain}
Let $\Omega$ be the domain \eqref{eq:omega-domain-intro}. It is fine-open
and invariant under $z\mapsto z+1$, $z\mapsto nz$, and
$z\mapsto z+k/n$ for the finite positive indices occurring in the Gamma
identities. For every assignment $a$,
\[
 \Lt_a(z)\in\No+i\Ocal\quad\Longleftrightarrow\quad z\in\Omega
 \qquad(z\in\tube).
\]
Thus
\begin{equation}\label{eq:complex-gamma}
 \Gt_a(z)=E(\Lt_a(z))\qquad(z\in\Omega)
\end{equation}
is a well-defined finely holomorphic, nowhere-zero surcomplex function.
It extends $\Gamma_a$ and satisfies
\begin{align}
 \Gt_a(z+1)&=z\Gt_a(z),\label{eq:actual-complex-recurrence}\\
 \Gt_a(nz)&=(2\pi)^{(1-n)/2}E((nz-\tfrac12)\log n)
 \prod_{k=0}^{n-1}\Gt_a(z+k/n).\label{eq:actual-complex-gauss}
\end{align}
\end{theorem}

\begin{proof}
At finite $z$ with nonzero standard part, the imaginary part of $\Lt_0(z)$
is a real constant plus an infinitesimal. At infinitesimal $z$, formula
\eqref{eq:complex-small} has finite imaginary part because $\operatorname{Arg}z$
is finite. The gauge has zero imaginary part and changes none of this.
At infinite real part,
\cref{lem:phase} gives exactly the condition defining $\Omega$.

To see openness at an infinite $x+iy\in\Omega$, restrict a perturbation
$u+iv$ to $\abs u<1$ and $\abs v<1/\log x$. The new real part is still
infinite and
$\log(x+u)-\log x=O(u/x)$. Therefore
$(y+v)\log(x+u)$ is finite. A further shrink preserves positivity.
The finite-real-part portion is already open. For the stated transformations,
$\log(nx)=\log x+\log n$ and
$\log(x+k/n)-\log x\in\mi$ at infinite $x$, while all imaginary
parts involved remain finite. The defining finite-product condition is
therefore preserved.

The complex derivative follows from the local chain rule for strong Taylor
expansions and $E'=E$. Nonvanishing follows from
$\abs{E(v+ib)}=e^v>0$. Apply the homomorphism $E$ to
\eqref{eq:complex-recurrence} and \eqref{eq:complex-gauss}; all arguments
have finite imaginary part, including $(nz-\tfrac12)\log n$.
Finally $E(\Log z)=z$, giving the claimed identities.
\end{proof}

\begin{example}[The boundary is a logarithmic scale]
At real part $x=\omega$, the point
$\omega+i/\log\omega$ lies in $\Omega$ and its log-Gamma phase is
$1$ plus an infinitesimal. The point $\omega+i$ lies in $\tube$ but not
in $\Omega$, because its log-Gamma phase is $\log\omega$ plus an
infinitesimal. On the other hand, a nonzero imaginary increment much smaller
than $1/\log\omega$ is permitted. Thus the obstruction is not simply
whether $y$ is finite or infinitesimal: it is the product $y\log x$.
\end{example}

\begin{warning}[An exact domain for one prescription, not a maximality claim]
Theorem \ref{thm:phase-domain} identifies exactly where the canonical
\emph{finite-phase exponentiation of this log-Gamma} is defined within
$\tube$. It does not prove that $\Gt_a$ has no other extension outside
$\Omega$, or that a global exponential cannot be chosen by additional
noncanonical phase data. Such choices define a different problem. No
value of $\sin\omega$ or $e^{i\omega}$ is selected in this article.
\end{warning}

\section{What the scalar derivation can distinguish}\label{sec:differential}

\subsection{The imported differential structure}
Let $\der$ be the Berarducci--Mantova derivation, normalized by
$\der\omega=1$. We use its Leibniz rule, strong additivity, exponential
compatibility, and constant field $\ker\der=\R$, all established in
\cite{BM}. In particular $\der\log x=\der x/x$ for $x>0$.
These are properties of scalar surreal numbers. No differentiation of
$M(x)$ or $\pp(x)$ as an ordinary coordinate function is assumed.

\begin{proposition}[The baseline satisfies the scalar chain rule]
\label{prop:BM-baseline}
For every positive surreal $x$,
\begin{equation}\label{eq:BM-chain0}
 \der(L_0(x))=L_0'(x)\der x.
\end{equation}
\end{proposition}
\begin{proof}
For a finite appreciable $x=r+\epsilon$, the coefficients of
\eqref{eq:L0-finite} are real constants for $\der$, and
$\der\epsilon=\der x$. Strong additivity and Leibniz give the
termwise Taylor chain rule. At positive infinitesimal $x$, the same argument
applies to the series in \eqref{eq:L0-small}, with the logarithmic term
handled by $\der\log x=\der x/x$. At infinite $x$, apply Leibniz and
logarithmic compatibility to the elementary terms in
\eqref{eq:L0-infinite}, and strong additivity to the Stirling series.
All differentiated families are summable by the strong-additivity property
of $\der$. The resulting expressions coincide with the fine derivatives
computed in \cref{prop:baseline-calculus}, multiplied by $\der x$.
\end{proof}

\subsection{Rigidity of infinitesimal locally constant discrepancies}
\begin{theorem}[Scalar compatibility removes flat local gauges]
\label{thm:BM-flat}
Let $h:\No_{>0}\to\No$ be finely locally constant, zero at positive
finite arguments, and infinitesimal at every positive infinite argument.
Set $L=L_0+h$. If
\begin{equation}\label{eq:BM-chain}
 \der(L(x))=L'(x)\der x\qquad(x>0),
\end{equation}
then $h=0$ and $L=L_0$.
In particular this conclusion applies to every algebraically flat gauge
constructed in this article.
\end{theorem}
\begin{proof}
Since $h'=0$, subtract \eqref{eq:BM-chain0} from
\eqref{eq:BM-chain}. It follows that $\der(h(x))=0$, so $h(x)\in\R$
for every $x$. The only real infinitesimal is zero. At infinite arguments
$h(x)$ therefore vanishes, and at finite ones it vanishes by hypothesis.
\end{proof}

The proof uses only infinitesimality of the discrepancy, not its full
algebraic flatness. The local-constancy hypothesis is essential to the stated
argument: we do not conclude uniqueness among arbitrary functions with a
scalar chain rule and a Stirling expansion.

\begin{theorem}[Complete scalar rigidity inside the monomialwise family]
\label{thm:BM-all}
For the family $L_a$ of \eqref{eq:gauge-intro}, the scalar chain rule
\eqref{eq:BM-chain} holds for all positive surreal arguments if and only if
$a=0$. No convexity or flatness assumption on $a$ is required.
\end{theorem}
\begin{proof}
The baseline case follows from \cref{prop:BM-baseline}. Conversely,
\cref{lem:gauge-identities} gives $L_a'=L_0'$, so the scalar chain rule
forces every $h_a(x)$ to be real. Fix an infinite monomial $m$ and put
$r=h_a(m)=m a(m)\in\R$. Since $\sqrt m$ is a smaller infinite monomial,
\[
 h_a(m+\sqrt m)=a(m)(m+\sqrt m)=r+\frac r{\sqrt m}.
\]
This too must be real. Thus $r/\sqrt m$ is real and infinitesimal, hence
zero. Therefore $r=0$ and $a(m)=0$. As $m$ was arbitrary, $a=0$.
\end{proof}

\begin{example}[An explicit derivative mismatch]
For $h_\lambda(x)$ from \eqref{eq:explicit-gauge}, its fine coordinate
derivative at $\omega$ is zero. In contrast, its scalar value there is
$\lambda\omega e^{-\omega}$, and
\begin{equation}\label{eq:BM-defect}
 \der(h_\lambda(\omega))
 =\lambda(1-\omega)e^{-\omega}.
\end{equation}
This is nonzero for $\lambda\ne0$. Consequently
\[
 \der(L_\lambda(\omega))-L_\lambda'(\omega)\der\omega
 =\lambda(1-\omega)e^{-\omega}.
\]
The discrepancy is explicit and exactly computed. It is not an assertion
that the two notions of derivative are inconsistent; they are different
operations and here they fail a particular compatibility requirement.
\end{example}

\subsection{Why scalar compatibility alone is not a universal uniqueness axiom}
The last theorem is an exact result for its named family. To make its scope
sharp, there are other, non-flat gauges which satisfy the scalar chain rule.

\begin{proposition}[A different real-valued gauge survives scalar compatibility]
\label{prop:real-gauge}
For $c_0\in\R$, define
\[
 k_{c_0}(x)=
 \begin{cases}
 0,&x\text{ positive finite},\\
 c_0\,\st\bigl(x/M(x)\bigr),&x\text{ positive infinite}.
 \end{cases}
\]
Then $L_0+k_{c_0}$ is strictly convex, satisfies the logarithmic recurrence
and every finite Gauss formula, has all the same positive-order fine
derivatives as $L_0$, and satisfies the scalar chain rule.
If $c_0\ne0$ it is not algebraically flat relative to $L_0$.
\end{proposition}
\begin{proof}
The function is real-valued, unchanged by finite translations, and scales
by $n$ under $x\mapsto nx$ for positive ordinary integers. Thus all the
identity and derivative arguments apply, and
$\der(k_{c_0}(x))=0$ because its value is real.
For convexity, finite or mixed arguments are treated as before. At infinite
arguments the gauge difference is finite. Whenever the baseline gap is
positive infinite this is harmless. In the remaining case $y/x\sim1$,
the leading monomial and leading real coefficient of $x$ and $y$ agree,
so $k_{c_0}(y)=k_{c_0}(x)$ exactly. Hence every tangent gap remains
positive. Finally $k_{c_0}(m)=c_0$ at an infinite monomial $m$, which is
not flat if $c_0\ne0$.
\end{proof}

Thus two distinct conclusions should not be confused. Scalar compatibility
eliminates every infinitesimal locally constant discrepancy, and it
eliminates every monomialwise-linear gauge $a(M(x))\pp(x)$.
It does not by itself eliminate every imaginable finite-valued locally
constant discrepancy. Flatness or an equally strong normalization has a real
role in a broader uniqueness statement.

\section{A reusable convexity principle}\label{sec:general-principle}
The proof is not restricted to the Bernoulli coefficients. It isolates an
abstract mechanism useful for other asymptotically convex special functions.
We state it with explicit hypotheses rather than claiming unproved
extensions to other functions.

\begin{theorem}[Scale-margin gauge principle]\label{thm:abstract}
Let $F:\No_{>0}\to\No$ be finely differentiable and have strictly
positive tangent gaps at distinct arguments. Suppose further that:
\begin{enumerate}[label=(\alph*)]
\item $D_F(y,x)$ is positive infinite whenever exactly one argument is
infinite, or when both are infinite and $y/x\not\sim1$ (including
infinite or infinitesimal ratios);
\item if $x\ne y$ are positive infinite and $y/x\sim1$, then
\[
 D_F(y,x)\sim\kappa\frac{(y-x)^2}{x}
\]
for a fixed positive real $\kappa$.
\end{enumerate}
For the gauges $h_a$ in \eqref{eq:gauge-intro},
$F+h_a$ is convex if and only if $m a(m)$ is finite for every
$m\in\Minf$. In the affirmative case it is strictly convex.
\end{theorem}
\begin{proof}
The proof of \cref{thm:sufficiency} uses exactly these hypotheses:
finite gauge differences are dominated by the gaps in (a), and in (b) the
ratio of a nonzero gauge difference to the baseline gap is asymptotic to
$\abs{a(m)}x/(\kappa\abs{y-x})$ when the increment is infinite.
It is infinitesimal under the proposed coefficient bound. For finite
increments the difference vanishes.
If the coefficient bound fails, use the same
$\delta=\pp\sqrt{\min(m,\abs{a(m)}m)}$ as in
\cref{thm:necessity}; the baseline gap is asymptotic to
$\kappa\delta^2/m$ and is dominated by the negative gauge increment.
This proves necessity and strict sufficiency.
\end{proof}

For $F=L_0$, \cref{lem:infinite-gaps,lem:mixed-gaps} verify all hypotheses
with $\kappa=1/2$. For another function, these hypotheses need separate
proofs, especially at extremely small relative increments. Knowing only
$F''(x)\sim2\kappa/x$ does not automatically establish them in the fine
surreal calculus. The abstract theorem states precisely what must be checked.

\section{Provenance, limitations, and further questions}\label{sec:provenance}

\subsection{Repository overlap audit}
The original repository inspection used the tree with observed commit
\begin{center}\small\commit\end{center}
and the documentation catalogue \cite{Repo}. That catalogue listed eighteen
research packages, including surreal construction and birthday questions,
surcomplex local and global analysis, polynomial algebra, trigonometry,
differential equations, rank-one analytic geometry, and spectral theory.
The differential-equations article was also inspected to avoid duplicating
its scalar-derivation and finite-phase results. This was a targeted catalogue
and relevant-source inspection, not a line-by-line audit of every archived
manuscript.

No other Gamma-specific research package was identified in that catalogue.
This is a historical account of that targeted inspection, not a current
exhaustive non-overlap claim. The maintained collection has since grown;
the original machine-readable audit in \texttt{data/source-audit.json}
is retained as a record of the earlier inspection.
The present contribution is therefore organized around Gamma uniqueness and
its sharp global convexity margin, rather than reproducing the repository's
existing general local-constancy, trigonometry, or differential-equation
results. The finite-phase exponential and scalar-derivation distinction are
background mechanisms; the log-Gamma phase calculation and the exact gauge
classification are the Gamma-specific results proved here.

\subsection{Published inputs and proposed contributions}
The imported components are kept separate from the claims of this article:
Conway normal forms and the usual surreal exponential
\cite{Conway,Gonshor}; positive-support Hahn substitution
\cite{Neumann}; restricted analytic and exponential transfer
\cite{vdDE}; classical Gamma identities, expansions, and remainder bounds
\cite{DLMF55,DLMF57,DLMF511,DLMF515}; the existing Taylor/transseries
extension approach of Costin--Ehrlich \cite{CE}; and the scalar derivation
of Berarducci--Mantova \cite{BM}.

The proposed contribution is the combination of the all-scale tangent gap
proof, the necessary-and-sufficient bound $m a(m)\in\Ocal$, the flat
counterexamples retaining all the stated identities and error inequalities,
the exact phase product criterion $y\log x\in\Ocal$, and the scalar
rigidity statements for the specified gauges. The short rigidity deductions
and abstract reformulation are consequences supporting the central theorem,
not separately advertised solutions of famous open problems.

The literature search included surreal Gamma extensions and
Bohr--Mollerup/non-Archimedean uniqueness terminology. The relevant
Costin--Ehrlich sections were read, and the classical facts were checked
against the NIST Digital Library of Mathematical Functions. I did not locate
this exact sharp coefficient theorem and combined invariance result in the
sources consulted. Search results for some exact phrases were sparse or
noisy. Accordingly, the justified status is \emph{proposed new results with
proofs}, not a certificate of priority or an exhaustive negative claim about
the entire literature.

\subsection{Claims not made}
No claim is made that the ordinary Gamma function has been newly constructed,
that the original real Bohr--Mollerup theorem is false, or that the
Costin--Ehrlich extension is internally ambiguous after its exact extension
operator has been fixed. Nor is the coefficient theorem a classification
of all surreal Gamma functions. We do not supply a proper-class contour
integral, a global surcomplex exponential, or a maximal surcomplex Gamma
domain. The asymptotic indices are ordinary integers, and the series are
strong sums, not limits of fine-topological partial sums.

The proofs have been algebraically checked in the finite computations
described in \cref{sec:verification}, but have not been formalized in Lean
or another proof assistant. In particular, computation does not replace the
all-scale order comparisons or the class-theoretic support arguments.
Independent verification of those proofs and of the scope of the novelty
claim remains appropriate before publication.

\subsection{Concrete next mathematical questions}
The exact family considered here suggests several focused problems, none of
which is claimed solved merely by the preceding results.

\paragraph{A full gauge classification.}
Describe all finite-translation-invariant $h$ for which $L_0+h$ is convex
and $h(nx)=nh(x)$ for every positive ordinary integer. The monomialwise
family and the real-valued leading-coefficient family show that there are
at least two genuinely different mechanisms. A complete answer would have
to describe allowable behavior on $\No/\Ocal$ rather than only on leading
monomials.

\paragraph{A minimal uniqueness package.}
Within locally constant infinitesimal discrepancies, scalar compatibility
is sufficient. Can it be replaced by a natural condition stated purely in
terms of the function and normal-form truncation? Any candidate must exclude
$\lambda\pp(x)e^{-M(x)}$ while preserving the baseline's local and
functional identities. Simply repeating all finite-order asymptotics or
all finite Gauss formulas cannot work.

\paragraph{Beyond the finite-imaginary tube.}
The log-Gamma construction and phase criterion here are restricted to
$\tube$. Extending a logarithmic prescription through much larger complex
sectors requires support control when both components are infinite. Actual
Gamma values then additionally require a specified phase structure. These
are two separate questions, not one unqualified request to exponentiate
an arbitrary surcomplex value.

\paragraph{Proof-assistant decomposition.}
A realistic formalization can isolate the ordered-field tangent criterion,
the algebraic gauge identities, and the abstract scale-margin theorem from
the analytic baseline. The first three do not require an implementation of
all surreal special functions. The remaining analytic obligations are
precisely the strong substitution lemmas, the finite restricted-analytic
transfer interface, and the four cases of the infinite tangent-gap lemma.
No uncompiled Lean code is presented as a verified formal proof.

\section{Conclusion}
The positive surreal half-line admits enough separation between finite
translations and infinite scales to preserve the entire local Gamma
calculus while changing the global function. Global log-convexity does not
remove that freedom: its exact tolerance in the monomialwise-linear family
is $a(m)=O(m^{-1})$. Exponentially small coefficients lie far inside the
permitted region, producing distinct normalized Gamma functions that share
all finite multiplication identities, all finite-translation normalized
germs, every Stirling coefficient, and every ordinary-order enveloping
Stirling inequality.

The phenomenon is not unrestricted arbitrariness. A coefficient larger than
the sharp threshold produces an explicit negative tangent gap. The
surcomplex logarithmic extension has a precise finite-phase condition. And
a scalar chain rule detects discrepancies invisible to fine coordinate
derivatives, removing all flat locally constant gauges in the specified
category. These three boundaries---convexity scale, phase scale, and
scalar compatibility---make the construction a controlled result rather
than an isolated pathological example.

\appendix
\section{Exact symbolic checks and reproduction}\label{sec:verification}

The archive includes \texttt{code/verify.py}, a Python program using SymPy,
and its recorded JSON output in \texttt{data/verification.json}.
The recorded run passed all 130 checks (Python 3.13.5 and SymPy 1.14.0;
the precise runtime versions are also saved in the JSON record).
The calculations use rational arithmetic and formal power series. There are
no floating-point tests pretending to represent infinite surreal numbers.

The tests check the Bernoulli normalizations, the shift cancellation
underlying \eqref{eq:formal-shift}, several finite Gauss defects, the tangent
kernel expansions, the quadratic divisibility of
\eqref{eq:R-gap-series}, and finite sparse normal-form identities for
$\pp$ and $M$. They also check the displayed derivative mismatch
\eqref{eq:BM-defect}. The normal-form tests are finite algebraic examples,
not a general implementation of the surreal field or its exponential.

The following commands reproduce the checks and build in temporary output
directories from the package directory, preserving the recorded evidence:
\begin{verbatim}
python -m pip install -r data/requirements.txt
python code/verify.py --output /tmp/gamma-verification-rerun.json
mkdir -p /tmp/gamma-article-build
latexmk -pdf -interaction=nonstopmode -halt-on-error \
  -outdir=/tmp/gamma-article-build article.tex
\end{verbatim}
Without \texttt{latexmk}, run \texttt{pdflatex} three times with
\texttt{-interaction=nonstopmode}, \texttt{-halt-on-error}, and
\path|-output-directory=/tmp/gamma-article-build|.

\Needspace{4\baselineskip}
The archived build helper is \texttt{code/makefile-Makefile}; its check target overwrites the historical
JSON record, so use the explicit temporary output path above for reruns.
The LaTeX source has an internal bibliography and no external graphics or
custom font files.

\subsection{A transparent recurrence calculation}
With $t=X^{-1}$ and $R_J(t)=\sum_{j=1}^Jb_jt^{2j-1}$, the formal recurrence
defect reduces to the ordinary formal expression
\begin{equation}\label{eq:verify-recurrence}
 \left(t^{-1}+\frac12\right)\log(1+t)-1
 +R_J\left(\frac{t}{1+t}\right)-R_J(t).
\end{equation}
Its coefficients through degree $2J+1$ vanish. The first coefficient that
can be affected by the omitted Stirling term has degree $2J+2$.
This finite truncation property is what the program checks. The full proof
of \cref{lem:formal-shift} establishes cancellation at every ordinary order.

\subsection{A transparent multiplication calculation}
For $u_k=k/n$, the elementary terms in the Gauss defect reduce, after
cancellation of $\log X$, $\log n$, and $c$, to
\[
 \sum_{k=0}^{n-1}\left[
 \left(t^{-1}+u_k-\frac12\right)\log(1+u_kt)-u_k\right].
\]
Adding
\[
 \sum_{k=0}^{n-1}R_J\left(\frac{t}{1+u_kt}\right)-R_J(t/n)
\]
gives the truncated defect tested by the program. Its coefficients through
degree $2J$ vanish; an omitted term of degree $2J+1$ can affect the next
coefficient. Checking more terms is a useful regression test, but the
periodic-series lemma is the proof at all orders.

\subsection{What the checks cannot establish}
The program does not decide positivity of arbitrary surreal functions,
quantify over a proper class, prove Neumann's support lemma, or check
first-order transfer. It also does not verify novelty. Its role is to expose
coefficient, sign, indexing, or normalization mistakes in the finite algebra
supporting the arguments. The actual proof obligations remain visible in
the numbered lemmas and theorems.

\begin{thebibliography}{99}

\bibitem{Conway}
J.~H. Conway,
\emph{On Numbers and Games}, second edition,
A.~K. Peters, 2001.

\bibitem{Gonshor}
H.~Gonshor,
\emph{An Introduction to the Theory of Surreal Numbers},
London Mathematical Society Lecture Note Series 110,
Cambridge University Press, 1986.

\bibitem{Neumann}
B.~H. Neumann,
On ordered division rings,
\emph{Transactions of the American Mathematical Society}
\textbf{66} (1949), 202--252.

\bibitem{vdDE}
L.~van den Dries and P.~Ehrlich,
Fields of surreal numbers and exponentiation,
\emph{Fundamenta Mathematicae} \textbf{167} (2001), 173--188;
erratum, \textbf{168} (2001), 295--297.
\href{https://doi.org/10.4064/fm167-2-3}{doi:10.4064/fm167-2-3}.

\bibitem{CE}
O.~Costin and P.~Ehrlich,
\emph{Integration on the surreals},
arXiv:2208.14331, version 5, July 4, 2024.
\url{https://arxiv.org/html/2208.14331v5}.
Especially the Taylor/transseries extension construction, \S10.1(v) on
Gamma and log-Gamma, and \S11 on set-sized fields and restricted analytic
transfer. The coefficient convention used here is specified independently
in \eqref{eq:coefficients}.

\bibitem{BM}
A.~Berarducci and V.~Mantova,
Surreal numbers, derivations and transseries,
\emph{Journal of the European Mathematical Society}
\textbf{20} (2018), 339--390.
\href{https://doi.org/10.4171/JEMS/769}{doi:10.4171/JEMS/769}.
\url{https://arxiv.org/abs/1503.00315}.

\bibitem{DLMF55}
NIST Digital Library of Mathematical Functions,
\S5.5, Functional Relations: recurrence, Gauss multiplication, and the
Bohr--Mollerup theorem.
\url{https://dlmf.nist.gov/5.5}.
Consulted September 2026.

\bibitem{DLMF57}
NIST Digital Library of Mathematical Functions,
\S5.7, Series Expansions.
\url{https://dlmf.nist.gov/5.7}.
Consulted September 2026.

\bibitem{DLMF59}
NIST Digital Library of Mathematical Functions,
\S5.9, Integral Representations, especially Binet's formula 5.9.10.
\url{https://dlmf.nist.gov/5.9}.
Consulted September 2026.

\bibitem{DLMF511}
NIST Digital Library of Mathematical Functions,
\S5.11, Asymptotic Expansions, including Eq.~5.11.1 and \S5.11(ii)
on signed error bounds.
\url{https://dlmf.nist.gov/5.11}.
Consulted September 2026.

\bibitem{DLMF515}
NIST Digital Library of Mathematical Functions,
\S5.15, Polygamma Functions.
\url{https://dlmf.nist.gov/5.15}.
Consulted September 2026.

\bibitem{Repo}
V.~Reshetnikov,
\emph{Surreal}, GitHub repository and documentation catalogue.
Observed tree revision \commit.
\url{https://github.com/VladimirReshetnikov/Surreal}.
Targeted inspection of \texttt{docs/README.md} and
\texttt{docs/surcomplex/differential-equations/article.tex}, September 2026.
This citation identifies the overlap audit, not a source for the new proofs.

\end{thebibliography}
\end{document}
